\documentclass[11pt]{article}

% ---------- packages ----------
\usepackage[margin=1in]{geometry}
\usepackage[T1]{fontenc}
\usepackage{lmodern}
\usepackage{microtype}
\usepackage{amsmath, amssymb, amsthm}
\usepackage{bm}
\usepackage{booktabs}
\usepackage{array}
\usepackage{tabularx}
\usepackage{longtable}
\usepackage{graphicx}
\usepackage{enumitem}
\usepackage{caption}
\usepackage{float}
\usepackage{xcolor}
\usepackage{titlesec}
\usepackage{fancyhdr}
\usepackage{listings}
\usepackage{tcolorbox}
\usepackage{hyperref}

% ---------- colors ----------
\definecolor{accent}{HTML}{1F4E79}
\definecolor{rule}{HTML}{2E75B6}
\definecolor{codebg}{HTML}{F4F4F4}
\definecolor{headfill}{HTML}{D9E2F3}
\definecolor{linkcolor}{HTML}{1F4E79}

\hypersetup{
  colorlinks=true,
  linkcolor=linkcolor,
  urlcolor=linkcolor,
  citecolor=linkcolor
}

% ---------- titles ----------
\titleformat{\section}
  {\color{accent}\Large\bfseries}{\thesection.}{0.6em}{}
  [{\color{rule}\vspace{-0.3em}\noindent\rule{\linewidth}{1pt}}]
\titleformat{\subsection}
  {\color{accent}\large\bfseries}{\thesubsection}{0.6em}{}
\titleformat{\subsubsection}
  {\bfseries}{\thesubsubsection}{0.6em}{}

\titlespacing*{\section}{0pt}{1.4em}{0.6em}
\titlespacing*{\subsection}{0pt}{1.1em}{0.4em}
\titlespacing*{\subsubsection}{0pt}{0.9em}{0.3em}

% ---------- headers / footers ----------
\pagestyle{fancy}
\fancyhf{}
\fancyhead[R]{\small\color{gray}Oyster Reef Restoration }
\fancyfoot[C]{\small\color{gray}Page \thepage\ of \pageref{LastPage}}
\renewcommand{\headrulewidth}{0.4pt}
\renewcommand{\headrule}{\hbox to\headwidth{\color{gray!50}\leaders\hrule height \headrulewidth\hfill}}

\usepackage{lastpage}

% ---------- spacing ----------
\setlength{\parskip}{0.6em}
\setlength{\parindent}{0pt}
\renewcommand{\arraystretch}{1.15}

% ---------- listings ----------
\lstdefinestyle{shell}{
  basicstyle=\ttfamily\footnotesize,
  backgroundcolor=\color{codebg},
  frame=single,
  framerule=0pt,
  framesep=6pt,
  xleftmargin=6pt,
  xrightmargin=6pt,
  breaklines=true,
  showstringspaces=false,
  columns=fullflexible
}
\lstset{style=shell}

% ---------- table column types ----------
\newcolumntype{S}{>{\ttfamily\footnotesize}l}    % monospace sites column
\newcolumntype{C}[1]{>{\centering\arraybackslash}p{#1}}
\newcolumntype{L}[1]{>{\raggedright\arraybackslash}p{#1}}

% ---------- helper boxes ----------
\newtcolorbox{notebox}{
  colback=codebg, colframe=rule, boxrule=0.4pt,
  arc=2pt, left=8pt, right=8pt, top=4pt, bottom=4pt
}

% ============================================================
\title{
  \vspace{-2em}
  \color{accent}\textbf{\Huge Oyster Reef Restoration Optimization}\\[0.4em]
  {\Large\color{black!70}\normalfont Site Selection under a JARS Metapopulation Connectivity Model}\\[0.3em]
  {\large\color{rule}\itshape Project Report}
}
\author{
  \textbf{Marouf Paul}\\[0.1em]
  \small M.S.\ Computational Operations Research, William \& Mary\\[0.1em]
  \small Advisor: Dr. Rex Kincaid \quad\textbullet\quad Collaborators: Dr.\ Leah Shaw, Dr.\ Rom Lipcius (VIMS)\\[0.1em]
  \small Repository: \texttt{NetworkAnalysisOysters}
}
\date{}

% ============================================================
\begin{document}
\maketitle
\thispagestyle{fancy}

% ============================================================
\section{Purpose of This Document}

This report has two jobs:

\begin{itemize}[leftmargin=1.4em, itemsep=0.2em]
  \item Describe the problem, the JARS metapopulation model, the optimization formulations, and the heuristic algorithms in enough detail that the code is reproducible without reading every script.
  \item Report only verified results. Every number in the tables below was re-derived from the repository by re-running the JARS ODE on the final selected sets, and every cross-model intersection was recomputed from the saved CSVs.
\end{itemize}

Throughout the report I use the following convention: greedy, post-greedy 1-swap, and stingy backward elimination are evaluated under the full JARS equilibrium objective $F(S)$, while all MIQP variants optimize a quadratic plumbing surrogate $\tilde{F}(S)$ using AMPL and Gurobi. Both views are needed because the surrogate is fast and provable but lossy, and the ODE is faithful but expensive. The surrogate retains self-recruitment (the connectivity diagonal), matching the JARS recruitment term, whose internal sum runs over all sites including the reef itself.

% ============================================================
\section{Problem Statement and Biology}

\subsection{Restoration as a network problem}

The Chesapeake Bay historically supported more than ten billion oysters; today the population is below $1\%$ of that baseline. Restoration is expensive, so the question is not whether to restore but \emph{where}: among $49$ candidate reef sites distributed across five ecological management zones in the bay, which subset of $K=25$ sites should be restored to maximize long-run adult oyster biomass?

The non-trivial part is that oyster larvae disperse via tidal currents. A reef's value depends not only on its own productivity but on the larvae it receives from, and sends to, the other selected reefs. This makes the problem combinatorial and non-additive: pairs and clusters of reefs interact super-linearly, and restoring an isolated high-productivity reef can be worse than restoring three weaker but well-connected ones.

\subsection{The JARS metapopulation ODE}

Each candidate reef is modeled by a coupled four-state ODE for juveniles $J$, adults $A$, reef shell volume $R$, and sediment $S$ (the JARS model, originally implemented by Prof.\ Leah Shaw and converted to Python in \texttt{src/model/jars\_ode.py}). The state of reef $k$ evolves as larval supply is filtered through density-dependent settlement, juveniles mature into adults, adults accrete shell, and sediment slowly buries shell. The key coupling term is the larval input
\[
\Lambda_k \;=\; P_0(k) \;+\; \sum_{\ell} P_1(\ell \to k)\,|A_\ell|^{1.72},
\]
where $P_0(k)$ is the external larval supply at site $k$, $P_1$ is the larval transport probability matrix between sites (the internal sum includes $\ell=k$, i.e.\ self-recruitment), and the exponent $1.72$ (a published parameter for oysters) makes dense, retentive clusters disproportionately valuable. This nonlinearity is the structural reason greedy heuristics undershoot and clustered selection methods do well.

\subsection{The optimization objective}

Let $A_k^*(S)$ be the equilibrium adult biomass at site $k$ when the ODE is integrated to $t=1000$ with reefs restored at the sites in $S$. The true objective is
\[
\max\;F(S) = \sum_{k \in S} A_k^*(S), \qquad \text{subject to } |S|=25,\; S \subseteq N,\; |N|=49.
\]
Every greedy / swap / stingy result reported below was scored by literally integrating this ODE on the chosen set. The MIQP models optimize a quadratic surrogate of $F$ that captures the dominant external + internal-retention contributions; the alignment between the two is excellent in practice but is not assumed in the heuristic results.

% ============================================================
\section{Connectivity Data and External Forcing}

\subsection{Two connectivity matrices}

VIMS collaborators provided two independent connectivity matrices from different hydrodynamic realizations:

\begin{itemize}[leftmargin=1.4em, itemsep=0.1em]
  \item \textbf{Matrix 1}: \texttt{nk\_All\_060102final\_56sites\_Model.xlsx}
  \item \textbf{Matrix 2}: \texttt{nk\_All\_060103final\_56sites\_Model.xlsx}
\end{itemize}

Both matrices include $56$ raw site labels. Per biological advice from Prof.\ Shaw and Prof.\ Lipcius, sites $66$--$72$ are excluded, leaving $49$ candidate sites for both matrices (\texttt{CANDIDATE\_SITES} in \texttt{src/model/jars\_ode.py}). For the MIQP, $P_1$ is additionally scaled by \texttt{P1scaling = 0.5} and fed through the surrogate weight $W(\ell, k) = P_1(\ell, k)\,(A^*)^{1.72}$ with $A^* = 0.05675$ representing the median single-reef adult equilibrium. The diagonal $W(k,k)$ is retained, so self-recruitment is credited in the surrogate exactly as it is in the ODE.

\subsection{Indexing convention (important)}

There are two coexisting indexing schemes in this project, and confusing them is the easiest way to get wrong results.

\begin{notebox}
\textbf{Biological site labels} are the original integer IDs in the Excel matrices: $1, 3, 4, 5, \ldots, 60$ (with sites $66$--$72$ dropped, and several others in the $1$--$60$ range absent). Every CSV file and every results table in this report uses these labels.

\textbf{AMPL indices} are $0, 1, \ldots, 48$, i.e.\ positions in the \texttt{CANDIDATE\_SITES} array. Every \texttt{.dat} file under \texttt{ampl/} uses these indices. The translation between the two lives in \texttt{runs/oyster\_index\_mapping\_matrix\{1,2\}.csv}.
\end{notebox}

In particular, the community sets in \texttt{ampl/oyster\_comm.dat} (e.g.\ $C_1 = \{33, 36, 44\}$) are AMPL \emph{indices}, not site labels. Looking up those indices in the mapping CSV shows that $C_1$ actually contains site labels $\{42, 48, 56\}$. The Python runner scripts (\texttt{scripts/run\_miqp*.py}) handle the translation back to site labels before writing CSVs, so end-user output is always in biological labels.

\subsection{External larval supply: two regimes}

The same set of $49$ sites is studied under two different external-forcing assumptions:

\begin{itemize}[leftmargin=1.4em, itemsep=0.2em]
  \item \textbf{Constant $P_0$} --- every site receives the same external input of $170$ larvae per timestep. This is the ``pure plumbing'' regime: differences in performance come only from the internal connectivity network, so this regime tells us about the structure of the network itself.
  \item \textbf{Realistic $P_0$} --- a heterogeneous vector with values in $\{100, 200, 400\}$ assigned per region based on empirical estimates (\texttt{setP0} in \texttt{jars\_ode.py}, supplied by the biologists). This regime mixes ``faucet strength'' into the optimization: some sites are inherently well-fed even if they are weakly connected.
\end{itemize}

All MIQP results in this report are under constant $P_0$ (the AMPL data file embeds $170$ for all sites). The realistic regime is only used for the ODE-driven heuristics.

% ============================================================
\section{Optimization Methods}

\subsection{ODE-driven heuristics}

\subsubsection*{Greedy forward selection}
Start with the empty set. At each step, evaluate $F(S \cup \{k\})$ by running the full JARS ODE for every candidate $k$ not yet selected, and add the site with the largest marginal gain. Repeat $25$ times. This is a fast first pass but myopic: it cannot recover from a poor early choice, and the $|A|^{1.72}$ nonlinearity means it tends to undershoot the best clustered design.

\subsubsection*{Post-greedy 1-for-1 local swap}
Starting from the greedy $25$-site set, evaluate every swap that exchanges one selected site for one unselected candidate. Apply the best improving swap, then repeat until no swap improves $F$. Empirically this lifts the greedy score by $1$--$2\%$ in absolute terms and is what we treat as the heuristic upper bound.

\subsubsection*{Stingy backward elimination}
Start with all $49$ candidates. At each step, evaluate $F(S \setminus \{a\})$ for every $a \in S$, and remove the site whose removal hurts $F$ the least. Stop when $|S|=25$. This is a separate path to the local optimum, and is useful as a consistency check on the swap result --- the two methods converging on the same $25$ sites is strong evidence that the set is a stable local optimum.

\subsection{MIQP surrogates}

All MIQP variants share the surrogate objective
\[
\max\;\tilde{F}(S) = \sum_{i \in S} Pe(i) + \sum_{i,j \in S} W(i, j)
\]
with binary selection variables $x_i$ and the cardinality constraint $\sum x_i = 25$. The double sum includes the diagonal term $W(i,i)$, which credits self-recruitment. There are four variants.

\subsubsection*{(a) Base MIQP --- \texttt{oyster\_quad.mod}}
Just the cardinality constraint. Picks the most retentive $25$-site cluster.

\subsubsection*{(b) MIQP + community minimums --- \texttt{oyster\_comm.mod}}
Adds five geographic-region constraints: $|S \cap C_r| \geq r_r$ with $(r_1, r_2, r_3, r_4, r_5) = (2, 5, 3, 2, 3)$. The total floor of $15$ leaves $10$ ``flex'' sites for the optimizer to place freely. This is the equity model used to satisfy multi-stakeholder coverage.

\subsubsection*{(c) MIQP + variable sizing --- \texttt{oyster\_size.mod}}
Adds a continuous reef-area variable $s_i$ for each candidate site, with
$L_i x_i \leq s_i \leq U_i x_i$, a total-area budget $\sum_i s_i \leq T=1000$ acres,
and a normalization $\bar{S}=20$. In the canonical size model used here, the
internal connectivity term uses \emph{source-only} area scaling:
\[
\max \; \sum_{i\in N} Pe(i)\,\frac{s_i}{\bar{S}}
      + \sum_{i,j\in N} W(i,j)\,\frac{s_i}{\bar{S}}\,x_j .
\]
Thus a larger source reef is credited with producing more larvae, while a
destination site only needs to be selected in order to receive internal larval
credit. This is the conservative size formulation: it does not multiply by the
recipient area $s_j$, avoiding possible double counting if recipient-site
footprint is already partly reflected in the raw fixed-release connectivity
matrix.

\subsubsection*{(d) MIQP + community + sizing --- \texttt{oyster\_comm\_size.mod}}
Combines (b) and (c). This is the model that most closely matches real planning constraints --- it must include some sites for regional equity, but it can allocate small areas to lower marginal-value sites while reserving the largest allowable areas for stronger larval source sites.

% ============================================================
\section{Results}

All ODE scores below were independently re-derived from the repository by running \texttt{src/opt/evaluator.py} on the saved site sets. All ten heuristic scores match the original logs exactly to six decimal places. MIQP objectives are read from the AMPL/Gurobi logs. The site-set intersections in \S5.3--5.4 were also recomputed.

\subsection{Heuristic results under constant \texorpdfstring{$P_0$}{P0}}

Table~\ref{tab:heur-const} lists the $25$-site sets and equilibrium scores produced by the three ODE-driven heuristics on each matrix under constant external supply.

\begin{table}[H]
\centering
\caption{ODE-driven heuristic results under constant $P_0$.}
\label{tab:heur-const}
\small
\begin{tabularx}{\linewidth}{@{}l X r@{}}
\toprule
\textbf{Matrix / Method} & \textbf{Selected sites (sorted)} & \textbf{$F(S)$} \\
\midrule
M1 Greedy & \texttt{4,\ 10,\ 15,\ 16,\ 17,\ 19,\ 20,\ 21,\ 26,\ 27,\ 28,\ 30,\ 31,\ 32,\ 36,\ 37,\ 40,\ 41,\ 44,\ 47,\ 49,\ 51,\ 52,\ 53,\ 54} & 1.846640 \\
M1 Greedy + Swap & \texttt{10,\ 12,\ 15,\ 16,\ 17,\ 20,\ 21,\ 26,\ 28,\ 29,\ 31,\ 32,\ 33,\ 36,\ 37,\ 40,\ 41,\ 44,\ 47,\ 49,\ 51,\ 52,\ 53,\ 54,\ 59} & 1.880054 \\
M1 Stingy & \texttt{10,\ 12,\ 15,\ 16,\ 17,\ 20,\ 21,\ 26,\ 28,\ 29,\ 31,\ 32,\ 33,\ 36,\ 37,\ 40,\ 41,\ 44,\ 47,\ 49,\ 51,\ 52,\ 53,\ 54,\ 59} & 1.880054 \\
\midrule
M2 Greedy & \texttt{4,\ 10,\ 11,\ 12,\ 15,\ 17,\ 19,\ 20,\ 21,\ 26,\ 27,\ 28,\ 29,\ 30,\ 31,\ 32,\ 36,\ 37,\ 40,\ 41,\ 47,\ 49,\ 51,\ 52,\ 53} & 1.692871 \\
M2 Greedy + Swap & \texttt{10,\ 11,\ 12,\ 15,\ 16,\ 17,\ 19,\ 20,\ 26,\ 29,\ 31,\ 32,\ 33,\ 36,\ 37,\ 40,\ 41,\ 44,\ 47,\ 49,\ 51,\ 52,\ 53,\ 54,\ 59} & 1.733033 \\
M2 Stingy & \texttt{10,\ 11,\ 12,\ 15,\ 16,\ 17,\ 19,\ 20,\ 26,\ 29,\ 31,\ 32,\ 33,\ 36,\ 37,\ 40,\ 41,\ 44,\ 47,\ 49,\ 51,\ 52,\ 53,\ 54,\ 59} & 1.733033 \\
\bottomrule
\end{tabularx}
\end{table}

Two structural observations from Table~\ref{tab:heur-const}:

\begin{itemize}[leftmargin=1.4em, itemsep=0.2em]
  \item On both matrices the swap improves on greedy by exchanging sites: on Matrix~1 it trades four sites ($\{4,19,27,30\}$ for $\{12,29,33,59\}$), raising $F$ from $1.846640$ to $1.880054$ ($+1.8\%$); on Matrix~2 it trades five sites ($\{4,21,27,28,30\}$ for $\{16,33,44,54,59\}$), raising $F$ from $1.692871$ to $1.733033$ ($+2.4\%$). The gain comes from changing the selected set, not from reordering a fixed one.
  \item On both matrices, post-greedy local swap and stingy backward elimination converge to exactly the same $25$ sites with exactly the same equilibrium score. This is strong (independent-path) evidence that these sets are stable local optima of the true ODE objective.
\end{itemize}

\subsection{Heuristic results under realistic \texorpdfstring{$P_0$}{P0}}

Switching to heterogeneous external supply pulls some new sites into the solution (those with high $P_0$ even if weakly connected) and pushes some plumbing-only sites out. Table~\ref{tab:heur-real} shows the greedy + swap and stingy results.

\begin{table}[H]
\centering
\caption{ODE-driven heuristic results under realistic $P_0$.}
\label{tab:heur-real}
\small
\begin{tabularx}{\linewidth}{@{}l X r@{}}
\toprule
\textbf{Matrix / Method} & \textbf{Selected sites (sorted)} & \textbf{$F(S)$} \\
\midrule
M1 Greedy + Swap (real) & \texttt{4,\ 6,\ 10,\ 15,\ 19,\ 20,\ 21,\ 24,\ 27,\ 30,\ 31,\ 32,\ 36,\ 37,\ 38,\ 39,\ 40,\ 41,\ 47,\ 49,\ 51,\ 52,\ 53,\ 55,\ 60} & 1.861968 \\
M1 Stingy (real) & \texttt{4,\ 10,\ 15,\ 16,\ 17,\ 19,\ 20,\ 21,\ 26,\ 27,\ 31,\ 32,\ 33,\ 36,\ 37,\ 40,\ 41,\ 44,\ 47,\ 49,\ 51,\ 52,\ 53,\ 54,\ 59} & 1.819074 \\
\midrule
M2 Greedy + Swap (real) & \texttt{1,\ 4,\ 6,\ 10,\ 18,\ 19,\ 20,\ 21,\ 24,\ 30,\ 31,\ 35,\ 36,\ 37,\ 38,\ 39,\ 40,\ 41,\ 42,\ 47,\ 49,\ 53,\ 55,\ 57,\ 60} & 1.793457 \\
M2 Stingy (real) & \texttt{4,\ 6,\ 7,\ 10,\ 15,\ 16,\ 19,\ 20,\ 21,\ 24,\ 26,\ 27,\ 30,\ 31,\ 32,\ 36,\ 37,\ 40,\ 41,\ 44,\ 47,\ 49,\ 51,\ 52,\ 53} & 1.723331 \\
\bottomrule
\end{tabularx}
\end{table}

Notice that under realistic forcing, greedy + swap and stingy no longer agree on a single set --- the heterogeneity of $P_0$ creates more local optima. However, a large ``robust core'' of sites is shared across all four realistic designs (see \S\ref{sec:backbone}).

\subsection{MIQP surrogate results under constant \texorpdfstring{$P_0$}{P0}}

All eight AMPL/Gurobi integer-programming models solved to proven global optimality with self-recruitment retained. Table~\ref{tab:miqp} reports the objective values and selected sets. Objective magnitudes differ between the baseline/community models and the reef-size models because the latter scale larval contributions by constructed reef area.

\begin{table}[H]
\centering
\caption{MIQP surrogate results under constant $P_0$ (self-recruitment retained, global optimum).}
\label{tab:miqp}
\small
\begin{tabularx}{\linewidth}{@{}l X r@{}}
\toprule
\textbf{Model} & \textbf{Selected sites (sorted)} & \textbf{Surrogate objective} \\
\midrule
M1 Base MIQP & \texttt{10,\ 11,\ 12,\ 15,\ 16,\ 17,\ 20,\ 26,\ 27,\ 28,\ 29,\ 31,\ 32,\ 33,\ 36,\ 37,\ 40,\ 41,\ 44,\ 49,\ 51,\ 52,\ 53,\ 54,\ 59} & 14785.03 \\
M1 MIQP + Comm & \texttt{1,\ 10,\ 12,\ 15,\ 16,\ 17,\ 21,\ 26,\ 28,\ 29,\ 31,\ 32,\ 33,\ 37,\ 40,\ 41,\ 42,\ 44,\ 47,\ 49,\ 51,\ 52,\ 53,\ 54,\ 59} & 13863.50 \\
M1 MIQP + Size & \texttt{10,\ 11,\ 12,\ 15,\ 16,\ 17,\ 21,\ 26,\ 27,\ 28,\ 29,\ 31,\ 32,\ 33,\ 36,\ 37,\ 40,\ 41,\ 44,\ 49,\ 51,\ 52,\ 53,\ 54,\ 59} & 32620.49 \\
M1 MIQP + Comm + Size & \texttt{1,\ 10,\ 12,\ 15,\ 16,\ 21,\ 26,\ 27,\ 28,\ 29,\ 31,\ 32,\ 33,\ 37,\ 40,\ 41,\ 42,\ 44,\ 47,\ 49,\ 51,\ 52,\ 53,\ 54,\ 59} & 30517.81 \\
\midrule
M2 Base MIQP & \texttt{10,\ 11,\ 12,\ 15,\ 16,\ 17,\ 20,\ 26,\ 27,\ 28,\ 29,\ 31,\ 32,\ 33,\ 36,\ 37,\ 40,\ 41,\ 44,\ 49,\ 51,\ 52,\ 53,\ 54,\ 59} & 16446.59 \\
M2 MIQP + Comm & \texttt{10,\ 11,\ 12,\ 15,\ 16,\ 17,\ 18,\ 19,\ 27,\ 28,\ 29,\ 31,\ 32,\ 33,\ 36,\ 37,\ 40,\ 41,\ 42,\ 47,\ 49,\ 51,\ 52,\ 53,\ 59} & 15294.18 \\
M2 MIQP + Size & \texttt{10,\ 11,\ 12,\ 15,\ 16,\ 17,\ 20,\ 26,\ 27,\ 28,\ 29,\ 31,\ 32,\ 33,\ 36,\ 37,\ 40,\ 41,\ 44,\ 49,\ 51,\ 52,\ 53,\ 54,\ 59} & 36695.67 \\
M2 MIQP + Comm + Size & \texttt{10,\ 11,\ 12,\ 15,\ 16,\ 17,\ 18,\ 19,\ 21,\ 27,\ 28,\ 29,\ 31,\ 32,\ 33,\ 36,\ 37,\ 40,\ 41,\ 42,\ 49,\ 51,\ 52,\ 53,\ 59} & 34226.75 \\
\bottomrule
\end{tabularx}
\end{table}

Two notes on the reef-sizing models. First, adding reef-size decisions changes the selected set only slightly relative to the corresponding model without reef-size decisions. The M2 size model selects the same $25$ sites as the M2 baseline MIQP, while the M1 size model exchanges one site (site $21$ for site $20$). Similarly, each community-plus-size model differs from its community-constrained counterpart by one site: site $27$ replaces site $17$ on M1, and site $21$ replaces site $47$ on M2. Within the selected set, the optimizer allocates the maximum allowable area ($40$--$50$ acres) to sites with the largest marginal source value, meaning sites whose external supply and outgoing larval contribution to the selected network are highest under the surrogate. Lower marginal-value selected sites are kept at the $5$-acre floor, with a single interior site absorbing the residual budget. For example, in the M1 MIQP + Size solution, the allocation is $50$ acres at sites $10, 15, 16, 17, 21, 26$--$29, 31$--$33$; $40$ acres at sites $36, 40, 41, 44, 49, 51$--$54$; $25$ acres at site $11$; and $5$ acres at sites $12, 37, 59$.

Second, the total reef area is exactly $1000$ acres in every sizing solution; the $1000$-acre constraint is tight in all cases.

\subsection{Backbone sites: intersections across models and matrices}
\label{sec:backbone}

The point of running several models is to see which sites are robust to the choice of method, connectivity matrix, and forcing regime. The sites in common across the solutions generated by each model provide a way to identify key reef site locations.

\subsubsection*{Within-matrix constant-$P_0$ intersection (all seven constant-forcing models)}

\begin{table}[H]
\centering
\small
\begin{tabular}{@{}lll@{}}
\toprule
\textbf{Matrix} & \textbf{Sites common to greedy, swap, stingy, MIQP, +Comm, +Size, +Comm+Size} & \textbf{Size} \\
\midrule
Matrix 1 & \texttt{10, 15, 16, 26, 28, 31, 32, 37, 40, 41, 44, 49, 51, 52, 53, 54} & 16 \\
Matrix 2 & \texttt{10, 11, 12, 15, 17, 29, 31, 32, 36, 37, 40, 41, 49, 51, 52, 53} & 16 \\
\bottomrule
\end{tabular}
\end{table}

This within-matrix table enumerates every constant-forcing run, including the greedy seed, applied symmetrically to both connectivity matrices. Matrix~1 has a $16$-site backbone and Matrix~2 a $16$-site backbone across the seven constant-forcing optimization runs.

\subsubsection*{Cross-matrix constant-$P_0$ backbone}

\begin{table}[H]
\centering
\small
\begin{tabular}{@{}lll@{}}
\toprule
\textbf{Selection} & \textbf{Sites} & \textbf{Size} \\
\midrule
Intersection of M1 (16) and M2 (16) backbones & \texttt{10, 15, 31, 32, 37, 40, 41, 49, 51, 52, 53} & 11 \\
\bottomrule
\end{tabular}
\end{table}

\subsubsection*{Global backbone (post-search and exact designs)}

\begin{table}[H]
\centering
\small
\begin{tabular}{@{}lll@{}}
\toprule
\textbf{Selection} & \textbf{Sites} & \textbf{Size} \\
\midrule
Sites selected by all $16$ post-search and exact designs & \texttt{10, 31, 37, 40, 41, 49, 53} & 7 \\
\bottomrule
\end{tabular}
\end{table}

The seven global-backbone sites --- \textbf{10, 31, 37, 40, 41, 49, 53} --- are selected by all $16$ post-search and exact designs considered here: the four swap heuristics and four stingy heuristics across both matrices and both external-forcing regimes, and the eight integer-programming solutions across both matrices and all four surrogate formulations. To avoid an asymmetric counting rule, standalone greedy is excluded in both forcing regimes, since greedy generates the initial solution for the local-search improvement procedure. The result does not depend on this choice. The two constant-forcing greedy seeds also contain the full backbone, and of the two realistic-forcing greedy seeds only the M1 seed omits a backbone site, site $37$. Counting only the post-search and exact designs gives sixteen of sixteen; counting the greedy seeds as well gives nineteen of twenty. Either way the backbone is the same set of seven sites.

\subsection{Network centrality on the global backbone}

To characterize how the seven backbone sites function in the transport network, we ranked all $49$ candidate reef restoration sites by five standard centrality measures: in-strength, out-strength, betweenness, eigenvector centrality, and PageRank. In- and out-strength are the off-diagonal weighted degree sums, and diagonal self-recruitment is reported separately. For betweenness, weights were converted to path lengths as $1/P_{i,k}$, so stronger larval-transfer links correspond to shorter paths. PageRank and eigenvector centrality were computed on the weighted directed graph with self-recruitment retained.

\begin{table}[H]
\centering
\caption{Network-centrality placement of the seven global-backbone sites among all $49$ candidates. Lower average rank is more central; rank $1$ is most central on a given metric.}
\label{tab:centrality}
\small
\begin{tabularx}{\linewidth}{@{}c c c X@{}}
\toprule
\textbf{Site} & \textbf{M1 avg. rank} & \textbf{M2 avg. rank} & \textbf{Most distinctive role} \\
\midrule
10 & 3.0  & 6.8  & strongest receiver (in-strength rank 1 on both matrices) \\
40 & 6.2  & 6.6  & spectral hub (eigenvector rank 1 on both matrices) \\
41 & 5.4  & 8.0  & strong source (out-strength rank 4--5) \\
31 & 8.0  & 5.6  & balanced hub, high on every metric \\
53 & 11.2 & 12.2 & connector (upper-tier betweenness) \\
49 & 15.8 & 10.2 & bridge (betweenness rank 1 on Matrix 2) \\
37 & 15.6 & 16.8 & complementary: modest standalone centrality \\
\bottomrule
\end{tabularx}
\end{table}

Four of the seven sites ($10,31,40,41$) are unambiguous hubs, sitting in the top eight of $49$ on average across both connectivity matrices. Sites $49$ and $53$ are valuable less as strength hubs than as connectors. Site $37$ is the informative exception: its standalone centrality is only moderate, and its out-strength is especially low, yet it is selected by every post-search and exact design. This illustrates the non-additivity that motivates the optimization problem. Site $37$ earns its place not by being individually central, but through its interaction with the rest of the selected set, which the quadratic surrogate rewards and a site-by-site centrality screen would miss.

% ============================================================
\section{Repository Structure and Reproducibility}

\subsection{Layout}

\begin{lstlisting}
NetworkAnalysisOysters/
|-- ampl/                   AMPL .mod and .dat files
|   |-- oyster_quad.mod     Base MIQP
|   |-- oyster_comm.mod     MIQP + community minimums
|   |-- oyster_size.mod     MIQP + variable sizing
|   |-- oyster_comm_size.mod  MIQP + community + sizing
|   `-- oyster_*.dat        Auto-generated data (prepare_miqp_data.py)
|-- data/                   Connectivity matrices (Excel)
|-- src/
|   |-- model/jars_ode.py   JARS ODE, candidate site list, P0 setter
|   |-- opt/evaluator.py    Single-call ODE scoring of a site set
|   |-- opt/greedy.py       Forward greedy
|   |-- opt/local_search.py 1-for-1 swap local search
|   |-- opt/backward.py     Stingy backward elimination
|   `-- utils/io_utils.py   CSV/TXT writers for runs/
|-- scripts/                Entry-point CLIs (run_*.py, prepare_*.py)
|-- runs/                   Saved CSVs, summaries, network metrics
|-- figures/                Generated PNGs
|-- validate.py             Independent re-validation harness
`-- requirements.txt
\end{lstlisting}

\subsection{Environment}

\begin{itemize}[leftmargin=1.4em, itemsep=0.1em]
  \item Python 3.10. \texttt{pip install -r requirements.txt} covers numpy, pandas, scipy, openpyxl, matplotlib, networkx, amplpy.
  \item AMPL + Gurobi on PATH. Test with \texttt{python -m scripts.check\_ampl\_env}.
\end{itemize}

\subsection{End-to-end reproduction}

All published numbers can be regenerated with the following sequence (substitute \texttt{--matrix 2} to repeat for the second connectivity matrix):

\begin{lstlisting}
# 1. Generate AMPL data from the Excel matrix (diagonal/self-recruitment kept)
python -m scripts.prepare_miqp_data --matrix 1

# 2. ODE-driven heuristics (constant P0)
python -m scripts.run_greedy
python -m scripts.run_greedy_then_local
python -m scripts.run_backward

# 3. MIQP variants (constant P0)
python -m scripts.run_miqp           --matrix 1
python -m scripts.run_miqp_comm      --matrix 1
python -m scripts.run_miqp_size      --matrix 1
python -m scripts.run_miqp_comm_size --matrix 1

# 4. Network centrality
python -m scripts.network_core_analysis --matrix 1

# 5. Figures
python -m scripts.make_figures
\end{lstlisting}

Run times on a recent laptop: each MIQP variant solves in well under a second; each greedy or stingy run takes $\sim 4$--$5$ minutes (because every step calls the ODE); a full greedy + local swap run takes $\sim 15$--$30$ minutes depending on how many improving passes are needed. A combined driver, \texttt{run\_everything.py}, parallelizes the per-candidate ODE evaluations across CPU cores and runs all heuristics and all eight MIQP solves in one pass; the parallel evaluation is bit-identical to the serial library functions.

\subsection{Validation script}

\texttt{validate.py} at the repository root is the independent re-validation harness for the handover. It reloads every reported design, re-runs the ODE on it, and recomputes all backbone intersections from scratch. Running it produces \texttt{runs/validation\_summary.csv} and \texttt{runs/validation\_network\_core\_matrix\{1,2\}.csv}. Every number in this report came from that script.

% ============================================================
\section{Limitations and Recommended Next Steps}

\subsection{Limitations of the current setup}

\begin{itemize}[leftmargin=1.4em, itemsep=0.2em]
  \item \textbf{Equilibrium-only objective.} $F(S)$ is the sum of $A$ at $t=1000$. Transient dynamics, time-to-recovery, and resilience to shocks are not considered.
  \item \textbf{Fixed JARS parameters.} The $14$ published JARS parameters are fixed; no robustness analysis was done over uncertainty in those parameters.
  \item \textbf{Surrogate weight uses a fixed $A^*$.} $W(i,j) = P_1(i,j)\,(0.05675)^{1.72}$ uses a single median single-reef equilibrium for $A^*$. A state-dependent surrogate could tighten the alignment between MIQP and ODE further, at the cost of solvability.
  \item \textbf{Community partition is exogenous.} The five communities come from geography, not from the connectivity graph. Data-driven community detection (Louvain, Infomap) on $P_1$ might produce more natural equity buckets.
\end{itemize}

\subsection{Self-recruitment sensitivity}

The canonical surrogate keeps self-recruitment (the connectivity diagonal). As a sensitivity check, zeroing the diagonal leaves the base and size \emph{selections} unchanged on both matrices and lowers the base objective by about $7\%$ on Matrix~1 ($14785 \to 13692$) and $4\%$ on Matrix~2 ($16447 \to 15807$). Only the community-constrained selections and the Matrix~1 within-matrix backbone shift; the cross-matrix and seven-site global backbones are unaffected. The headline result is therefore robust to whether self-recruitment is credited in the surrogate, while the canonical model keeps it for consistency with the ODE.

\subsection{Recommended next experiments}

\begin{itemize}[leftmargin=1.4em, itemsep=0.2em]
  \item \textbf{Robust / stochastic MIQP.} Treat the connectivity matrix entries as uncertain and solve a chance-constrained or two-stage stochastic version. The handful of disagreements between Matrix~1 and Matrix~2 are exactly the sites where this would matter.
  \item \textbf{Multi-objective frontier.} Trace the Pareto frontier of (biomass) vs.\ (equity tightness) vs.\ (total acreage). The current four models are four points on that frontier; ten or twenty would tell stakeholders what they're really trading.
  \item \textbf{Phased restoration.} Replace the one-shot $|S|=25$ selection with a multi-period plan that respects an annual restoration capacity. This would change the optimal sequence even if the final set stays similar.
  \item \textbf{Coupled with non-larval services.} The current objective is adult biomass. A more honest objective would also credit water filtration, fish habitat, and economic value at each site, weighted by stakeholder preferences.
\end{itemize}

\subsection{Small cleanups for whoever takes over}

\begin{itemize}[leftmargin=1.4em, itemsep=0.2em]
  \item \texttt{scripts/network\_core\_analysis.py} hard-codes a core-site list at the top. The validated backbones from this re-analysis are: global $7$ sites (\texttt{10, 31, 37, 40, 41, 49, 53}), cross-matrix $11$ sites, and within-matrix $16$ sites on both M1 and M2. The hard-coded list should be regenerated from these validated intersections.
  \item \texttt{scripts/prepare\_miqp\_data.py} contains a commented-out \texttt{np.fill\_diagonal(P, 0.0)} with a ``remove self-loops'' comment. Self-recruitment is intentionally retained, so that line should stay disabled and the comment reworded to avoid confusion.
  \item \texttt{scripts/experimental/run\_miqp\_inline.py} and \texttt{run\_miqp\_test.py} \emph{do} zero the diagonal. They are not part of the canonical pipeline; either delete the \texttt{experimental/} folder or comment those lines so the diagonal-zeroed surrogate is never run by accident.
  \item \texttt{ampl/oyster\_comm.dat} uses integer indices $0$--$48$ (the position in \texttt{CANDIDATE\_SITES}, not the original site label). Worth a one-line comment in the file since the same \texttt{.dat} is read alongside \texttt{.mod} files where $N = \{0, \ldots, 48\}$.
  \item \texttt{runs/} contains $\sim 60$ timestamped greedy CSVs from the development iterations. A \texttt{runs/final/} subfolder with just the canonical CSVs would make the repo friendlier on first read.
  \item A LibreOffice / Office lock file \texttt{\~{}\$presentation.pptx} is checked in. Safe to delete.
\end{itemize}

% ============================================================
\section{Conclusions}

Across two contrasting connectivity matrices, two external-forcing regimes, ODE-driven post-search designs, and four integer-programming variants, the same seven reefs are selected by every post-search and exact design: $10, 31, 37, 40, 41, 49, 53$. The eleven-site cross-matrix constant-forcing backbone $\{10, 15, 31, 32, 37, 40, 41, 49, 51, 52, 53\}$ is the next natural layer out, and the per-matrix backbones (sixteen sites on Matrix~1 and sixteen on Matrix~2) extend that further.

That this much agreement emerges from objectives that are conceptually different --- the full nonlinear JARS equilibrium on one hand, a static quadratic surrogate on the other --- is not a tautology. It tells us that the surrogate captures the dominant structural drivers of the metapopulation dynamics for this instance. From a planning perspective, the practical translation is that a small number of reefs do much of the connectivity work, and the seven-site backbone identifies high-priority candidates within the modeled connectivity and demographic assumptions.

% ============================================================
\appendix

% ============================================================
\section{All 18 Final Designs}

For ease of cross-checking, every $25$-site set used in this report is reproduced below as a flat list.

\begin{longtable}{@{}l p{12cm}@{}}
\toprule
\textbf{Model} & \textbf{Selected sites (sorted)} \\
\midrule
\endfirsthead
\toprule
\textbf{Model} & \textbf{Selected sites (sorted)} \\
\midrule
\endhead
M1 Greedy (constant) & \texttt{4, 10, 15, 16, 17, 19, 20, 21, 26, 27, 28, 30, 31, 32, 36, 37, 40, 41, 44, 47, 49, 51, 52, 53, 54} \\
M1 Greedy + Swap (constant) & \texttt{10, 12, 15, 16, 17, 20, 21, 26, 28, 29, 31, 32, 33, 36, 37, 40, 41, 44, 47, 49, 51, 52, 53, 54, 59} \\
M1 Stingy (constant) & \texttt{10, 12, 15, 16, 17, 20, 21, 26, 28, 29, 31, 32, 33, 36, 37, 40, 41, 44, 47, 49, 51, 52, 53, 54, 59} \\
M1 MIQP (constant) & \texttt{10, 11, 12, 15, 16, 17, 20, 26, 27, 28, 29, 31, 32, 33, 36, 37, 40, 41, 44, 49, 51, 52, 53, 54, 59} \\
M1 MIQP + Comm & \texttt{1, 10, 12, 15, 16, 17, 21, 26, 28, 29, 31, 32, 33, 37, 40, 41, 42, 44, 47, 49, 51, 52, 53, 54, 59} \\
M1 MIQP + Size & \texttt{10, 11, 12, 15, 16, 17, 21, 26, 27, 28, 29, 31, 32, 33, 36, 37, 40, 41, 44, 49, 51, 52, 53, 54, 59} \\
M1 MIQP + Comm + Size & \texttt{1, 10, 12, 15, 16, 21, 26, 27, 28, 29, 31, 32, 33, 37, 40, 41, 42, 44, 47, 49, 51, 52, 53, 54, 59} \\
M1 Greedy + Swap (realistic) & \texttt{4, 6, 10, 15, 19, 20, 21, 24, 27, 30, 31, 32, 36, 37, 38, 39, 40, 41, 47, 49, 51, 52, 53, 55, 60} \\
M1 Stingy (realistic) & \texttt{4, 10, 15, 16, 17, 19, 20, 21, 26, 27, 31, 32, 33, 36, 37, 40, 41, 44, 47, 49, 51, 52, 53, 54, 59} \\
M2 Greedy (constant) & \texttt{4, 10, 11, 12, 15, 17, 19, 20, 21, 26, 27, 28, 29, 30, 31, 32, 36, 37, 40, 41, 47, 49, 51, 52, 53} \\
M2 Greedy + Swap (constant) & \texttt{10, 11, 12, 15, 16, 17, 19, 20, 26, 29, 31, 32, 33, 36, 37, 40, 41, 44, 47, 49, 51, 52, 53, 54, 59} \\
M2 Stingy (constant) & \texttt{10, 11, 12, 15, 16, 17, 19, 20, 26, 29, 31, 32, 33, 36, 37, 40, 41, 44, 47, 49, 51, 52, 53, 54, 59} \\
M2 MIQP (constant) & \texttt{10, 11, 12, 15, 16, 17, 20, 26, 27, 28, 29, 31, 32, 33, 36, 37, 40, 41, 44, 49, 51, 52, 53, 54, 59} \\
M2 MIQP + Comm & \texttt{10, 11, 12, 15, 16, 17, 18, 19, 27, 28, 29, 31, 32, 33, 36, 37, 40, 41, 42, 47, 49, 51, 52, 53, 59} \\
M2 MIQP + Size & \texttt{10, 11, 12, 15, 16, 17, 20, 26, 27, 28, 29, 31, 32, 33, 36, 37, 40, 41, 44, 49, 51, 52, 53, 54, 59} \\
M2 MIQP + Comm + Size & \texttt{10, 11, 12, 15, 16, 17, 18, 19, 21, 27, 28, 29, 31, 32, 33, 36, 37, 40, 41, 42, 49, 51, 52, 53, 59} \\
M2 Greedy + Swap (realistic) & \texttt{1, 4, 6, 10, 18, 19, 20, 21, 24, 30, 31, 35, 36, 37, 38, 39, 40, 41, 42, 47, 49, 53, 55, 57, 60} \\
M2 Stingy (realistic) & \texttt{4, 6, 7, 10, 15, 16, 19, 20, 21, 24, 26, 27, 30, 31, 32, 36, 37, 40, 41, 44, 47, 49, 51, 52, 53} \\
\bottomrule
\end{longtable}

% ============================================================
\section*{Contact}

\noindent\textbf{Marouf Paul} --- \texttt{mmarouf@wm.edu} / \texttt{maroufpaul2@gmail.com} (post-graduation)

\noindent Advisor: Prof.\ Rex Kincaid (William \& Mary, Mathematics)

\noindent Biological collaborators: Prof.\ Leah Shaw (W\&M), Prof.\ Rom Lipcius (VIMS)

\end{document}
